% from ega3-7-fr.tex
\oldpage[III]{209}
\[
  y\longmapsto
  \sum_i(-1)^i\operatorname{rang}_{\kappa(y)}
  \mathscr T_i\bigl(\kappa(y)\bigr)
  \tag{7.9.4.1}\label{III.7.9.4.1.fr}
\]
est localement constante dans $Y$.
\end{theorem}

\begin{proof}
On peut se borner au cas où $Y=\operatorname{Spec}(A)$ est affine d'anneau $A$
noethérien. Comme le complexe $\mathscr P_\bullet$ est fini, on sait (7.7.12,
(i)) qu'on a
\[
  \mathscr T_p(\mathscr M)
  =\mathscr H_p(\mathscr L_\bullet\otimes_{\mathscr O_Y}\mathscr M),
  \qquad \mathscr L_\bullet=\widetilde{L_\bullet},
\]
$L_\bullet$ étant un complexe \emph{fini} de $A$-modules projectifs de type
fini. Le théorème résulte alors de (7.9.3).
\end{proof}

\begin{env}[7.9.5]
\label{III.7.9.5.fr}
Sous les conditions de (7.9.4), la fonction (7.9.4.1) est \emph{constante}
lorsque $Y$ est \emph{connexe}. Lorsque $Y$ est connexe et non vide, on désigne
la valeur unique (entière) de (7.9.4.1) par
\[
  \operatorname{EP}(f,\mathscr P_\bullet)
  \quad\text{ou}\quad
  \operatorname{EP}(Y,\mathscr P_\bullet),
  \quad\text{ou simplement}\quad
  \operatorname{EP}(\mathscr P_\bullet)
\]
s'il ne peut en résulter de confusion, et l'on dit que cet entier est la
\emph{caractéristique d'Euler-Poincaré de $\mathscr P_\bullet$} relativement à
$f$ (ou à $Y$). Dans le cas général, on notera aussi
\[
  \operatorname{EP}(f,\mathscr P_\bullet;y)
  \quad\text{ou}\quad
  \operatorname{EP}(Y,\mathscr P_\bullet;y)
  \quad\text{ou}\quad
  \operatorname{EP}(\mathscr P_\bullet;y)
\]
le second membre de (7.9.4.1).
\end{env}

\begin{env}[7.9.6]
\label{III.7.9.6.fr}
Sous les hypothèses de (7.9.4) relativement à $X$, $Y$ et $f$, soit
\[
  0\longrightarrow\mathscr P'_\bullet
  \xrightarrow{u}\mathscr P_\bullet
  \xrightarrow{v}\mathscr P''_\bullet
  \longrightarrow0
\]
une suite exacte de complexes finis de $\mathscr O_X$-Modules cohérents et
$Y$-plats, les homomorphismes $u$ et $v$ étant de degrés \emph{pairs} $2d$,
$2d'$ respectivement. Comme $\mathscr T_\bullet$ est un foncteur homologique
(7.7.1), on a une suite exacte d'homologie
\[
  \cdots\longrightarrow
  \mathscr T_i(\mathscr P'_\bullet,\kappa(y))
  \longrightarrow
  \mathscr T_{i+2d}(\mathscr P_\bullet,\kappa(y))
  \longrightarrow
  \mathscr T_{i+2d+2d'}(\mathscr P''_\bullet,\kappa(y))
  \longrightarrow
  \mathscr T_{i-1}(\mathscr P'_\bullet,\kappa(y))
  \longrightarrow\cdots
\]
n'ayant d'ailleurs qu'un nombre fini de termes. En écrivant que la
caractéristique d'Euler-Poincaré de ce complexe est nulle ($0_{\mathrm I}$,
11.10.1), il vient aussitôt
\[
  \operatorname{EP}(\mathscr P_\bullet;y)
  =\operatorname{EP}(\mathscr P'_\bullet;y)
  +\operatorname{EP}(\mathscr P''_\bullet;y)
  \tag{7.9.6.1}\label{III.7.9.6.1.fr}
\]
pour tout $y\in Y$. Or, si par exemple $\mathscr P_\bullet=(\mathscr P_i)$ avec
$\mathscr P_i=0$ pour $i<0$, on a la suite exacte de complexes
\[
\xymatrix@C=2.4em{
  \cdots \ar[r] & 0 \ar[r]\ar[d] & 0 \ar[r]\ar[d]
    & 0 \ar[r]\ar[d] & 0 \ar[r] & \cdots \\
  \cdots \ar[r] & 0 \ar[r]\ar[d] & 0 \ar[r]\ar[d]
    & \mathscr P_1 \ar[r]\ar[d] & \mathscr P_2 \ar[r] & \cdots \\
  \cdots \ar[r] & 0 \ar[r]\ar[d] & \mathscr P_0 \ar[r]\ar[d]
    & \mathscr P_1 \ar[r]\ar[d] & \mathscr P_2 \ar[r] & \cdots \\
  \cdots \ar[r] & 0 \ar[r]\ar[d] & \mathscr P_0 \ar[r]\ar[d]
    & 0 \ar[r]\ar[d] & 0 \ar[r] & \cdots \\
  \cdots \ar[r] & 0 \ar[r] & 0 \ar[r]
    & 0 \ar[r] & 0 \ar[r] & \cdots
}
\]
les flèches verticales non nulles étant les automorphismes identiques ; on peut
appliquer (7.9.6.1) à cette suite exacte, d'où, par récurrence sur la longueur de
$\mathscr P_\bullet$, la formule
\[
  \operatorname{EP}(\mathscr P_\bullet;y)
  =\sum_i(-1)^i\operatorname{EP}(\mathscr P_i;y).
  \tag{7.9.6.2}\label{III.7.9.6.2.fr}
\]

